% =====================================================================
%  GraficosEZ - Analizador de CSV
%  Manual de usuario
%
%  Compilar (dos pasadas para el índice):
%     pdflatex manual.tex && pdflatex manual.tex
%  o bien:
%     latexmk -pdf manual.tex
%  También compila en Overleaf (subir manual.tex y la carpeta img/).
% =====================================================================
\documentclass[11pt,a4paper,oneside]{report}

% ---------- idioma y tipografía ----------
\usepackage[utf8]{inputenc}
\usepackage[T1]{fontenc}
\usepackage{lmodern}
\usepackage[spanish,es-tabla,es-noshorthands,es-nodecimaldot]{babel}
\usepackage{microtype}

% ---------- página ----------
\usepackage[margin=2.5cm,headheight=14pt]{geometry}
\usepackage{fancyhdr}
\usepackage{titlesec}
\usepackage{parskip}

% ---------- elementos ----------
\usepackage{xcolor}
\usepackage{graphicx}
\usepackage{booktabs}
\usepackage{tabularx}
\usepackage{array}
\usepackage{enumitem}
\usepackage{listings}
\usepackage{tcolorbox}
\usepackage{amsmath,amssymb}
\usepackage[hidelinks]{hyperref}

\graphicspath{{img/}{../../src/assets/}}

% ---------- colores (tomados de la interfaz) ----------
\definecolor{acento}{HTML}{3B6EF6}
\definecolor{tinta}{HTML}{0F172A}
\definecolor{suave}{HTML}{64748B}
\definecolor{linea}{HTML}{E4E7EC}
\definecolor{fondo}{HTML}{F3F4F7}
\definecolor{okc}{HTML}{15803D}
\definecolor{warnc}{HTML}{B45309}
\definecolor{critc}{HTML}{B91C1C}

\hypersetup{colorlinks=true,linkcolor=acento,urlcolor=acento,
  pdftitle={GraficosEZ - Manual de usuario},pdfauthor={Tomas Sandoval}}

% ---------- títulos ----------
\titleformat{\chapter}[display]
  {\normalfont\sffamily\bfseries\color{tinta}}
  {\large\color{acento}Capítulo \thechapter}{4pt}{\Huge}
\titlespacing*{\chapter}{0pt}{-10pt}{24pt}
\titleformat*{\section}{\normalfont\sffamily\Large\bfseries\color{tinta}}
\titleformat*{\subsection}{\normalfont\sffamily\large\bfseries\color{tinta}}

% ---------- encabezado y pie ----------
\pagestyle{fancy}
\fancyhf{}
\fancyhead[L]{\small\sffamily\color{suave}GraficosEZ · Manual de usuario}
\fancyhead[R]{\small\sffamily\color{suave}\leftmark}
\fancyfoot[C]{\small\sffamily\thepage}
\renewcommand{\headrulewidth}{0.4pt}
\fancypagestyle{plain}{\fancyhf{}\fancyfoot[C]{\small\sffamily\thepage}\renewcommand{\headrulewidth}{0pt}}

% ---------- listas ----------
\setlist{itemsep=2pt,topsep=4pt}

% ---------- código ----------
\lstdefinestyle{json}{
  basicstyle=\ttfamily\small,
  backgroundcolor=\color{fondo},
  frame=single, rulecolor=\color{linea},
  framesep=6pt, xleftmargin=8pt, xrightmargin=8pt,
  breaklines=true, columns=fullflexible, keepspaces=true,
  showstringspaces=false,
  literate={á}{{\'a}}1 {é}{{\'e}}1 {í}{{\'i}}1 {ó}{{\'o}}1 {ú}{{\'u}}1
           {Á}{{\'A}}1 {É}{{\'E}}1 {Í}{{\'I}}1 {Ó}{{\'O}}1 {Ú}{{\'U}}1
           {ñ}{{\~n}}1 {Ñ}{{\~N}}1 {…}{{\ldots}}3 {·}{{\textperiodcentered}}1
}
\lstset{style=json}

% ---------- cajas ----------
\newtcolorbox{nota}[1][Nota]{colback=acento!5,colframe=acento,
  fonttitle=\sffamily\bfseries,title=#1,arc=2mm,boxrule=0.6pt}
\newtcolorbox{aviso}[1][Importante]{colback=warnc!6,colframe=warnc,
  fonttitle=\sffamily\bfseries,title=#1,arc=2mm,boxrule=0.6pt}
\newtcolorbox{consejo}[1][Consejo]{colback=okc!5,colframe=okc,
  fonttitle=\sffamily\bfseries,title=#1,arc=2mm,boxrule=0.6pt}

% ---------- macros ----------
\newcommand{\app}{\textsf{\textbf{GraficosEZ}}}
\newcommand{\boton}[1]{\fcolorbox{linea}{white}{\small\sffamily #1}}
\newcommand{\pestana}[1]{\textsf{\textbf{\color{acento}#1}}}
\newcommand{\campo}[1]{\textsf{\textit{#1}}}
\newcommand{\bs}{\textbackslash}
\newcommand{\warnz}{\textsf{\textbf{\color{warnc}WARNING}}}
\newcommand{\critz}{\textsf{\textbf{\color{critc}CRITICAL}}}
\newcommand{\okz}{\textsf{\textbf{\color{okc}OK}}}
\newcolumntype{L}[1]{>{\raggedright\arraybackslash}p{#1}}

% =====================================================================
\begin{document}

% ---------------------------- portada ----------------------------
\begin{titlepage}
  \centering
  \vspace*{2.5cm}
  \includegraphics[width=0.55\textwidth]{logo-app-light.png}\par
  \vspace{1.5cm}
  {\sffamily\Huge\bfseries\color{tinta} Manual de usuario\par}
  \vspace{0.4cm}
  {\sffamily\Large\color{suave} Analizador de CSV con series de tiempo\\
   y alertas por umbrales\par}
  \vspace{2cm}
  {\color{acento}\rule{0.6\textwidth}{1pt}}\par
  \vspace{0.6cm}
  {\sffamily\large Versión 1.0\par}
  \vspace{0.2cm}
  {\sffamily\normalsize\color{suave}\url{https://tomsandovalmed.github.io/GraficosEZ/}\par}
  \vfill
  \vspace{0.3cm}
  {\sffamily\color{suave} Desarrollado por Tomas Sandoval\par}
  \vspace{0.2cm}
  {\sffamily\small\color{suave} Noviembre de 2026\par}
\end{titlepage}

\tableofcontents

% =====================================================================
\chapter{Introducción}

\app{} (también llamado \emph{Analizador de CSV}) es una herramienta web para
\textbf{explorar, comparar y auditar series de tiempo} contenidas en archivos
CSV. Fue pensada especialmente para los registros de rendimiento de
\emph{Performance Monitor} (Perfmon) de Windows y SQL Server, pero funciona con
cualquier CSV que tenga una columna de fecha/hora y columnas numéricas.

\section{Características principales}

\begin{itemize}
  \item \textbf{Funciona 100\,\% en el navegador.} Los archivos se leen en tu
        propio equipo; no se suben a ningún servidor.
  \item \textbf{Lectura automática de los CSV:} detecta el separador, la
        columna de fecha, el formato de fecha, las columnas numéricas y el
        servidor de origen, sin configurar nada.
  \item \textbf{Reconocimiento de Perfmon:} entiende cabeceras del tipo
        \texttt{\bs\bs SERVIDOR\bs Objeto\bs Contador}, incluidas las instancias
        con nombre de SQL Server.
  \item \textbf{Gráficos interactivos:} un gráfico por contador, varios
        servidores superpuestos, zoom con el ratón, lectura de valores al pasar
        el cursor y franjas de color \warnz{} y \critz.
  \item \textbf{Alertas automáticas:} agrupa las muestras fuera de rango en
        episodios con fecha, severidad, valor pico y duración, y permite saltar
        a cada uno en el gráfico.
  \item \textbf{Comparación de periodos:} superpone dos rangos de fechas (del
        mismo archivo o de archivos distintos) y calcula la variación
        porcentual.
  \item \textbf{Perfiles de umbrales editables:} cada empresa o cliente tiene su
        propio perfil; se crean, editan, importan y exportan como archivos
        \texttt{.json}. Incluye los perfiles de fábrica \textbf{SII} y
        \textbf{TGR}.
  \item \textbf{Exportación de gráficos} a PNG o JPEG, o copia directa al
        portapapeles, con la tabla de estadísticas incluida en la imagen.
  \item \textbf{Diseño adaptable:} se usa en escritorio, tableta o teléfono, y
        respeta el modo claro u oscuro del sistema.
\end{itemize}

\section{Privacidad}

Todo el procesamiento ocurre en la pestaña del navegador. Los archivos CSV no
se envían a Internet. Lo único que la aplicación guarda es la lista de
\textbf{perfiles de umbrales}, en el almacenamiento local del navegador
(\emph{localStorage}). Los archivos cargados \textbf{no} se guardan: al
recargar la página hay que volver a cargarlos.

\section{Convenciones de este manual}

\begin{itemize}
  \item \pestana{Explorar}: nombre de una pestaña de la aplicación.
  \item \boton{Descargar gráfico}: un botón o control que se pulsa.
  \item \campo{Valor ideal}: un campo o columna de una tabla.
\end{itemize}

% =====================================================================
\chapter{Acceso y requisitos}

\section{Abrir la aplicación}

Hay dos formas de usarla:

\begin{enumerate}
  \item \textbf{En línea (recomendado).} Abre
        \url{https://tomsandovalmed.github.io/GraficosEZ/} en el navegador. La
        página se actualiza sola cada vez que se publica una nueva versión.
  \item \textbf{Sin conexión.} El proyecto se compila en un único archivo,
        \texttt{dist/index.html}, sin dependencias externas. Basta con
        descargarlo y abrirlo con doble clic. También se puede copiar a una
        carpeta compartida o a cualquier hosting estático.
\end{enumerate}

\section{Navegadores compatibles}

Cualquier navegador moderno y actualizado (Chrome, Edge, Firefox o Safari).

\begin{nota}
La función \boton{Copiar al portapapeles} para imágenes depende de que el
navegador lo permita. Algunos navegadores solo lo admiten en páginas
\texttt{https} (como la versión en línea) y pueden bloquearlo al abrir el
archivo local. Si falla, la aplicación lo avisa; usa entonces
\boton{Descargar gráfico}.
\end{nota}

\section{Versión}

La versión instalada aparece en la esquina inferior derecha de la página
(por ejemplo, \emph{Versión 1.0}).

% =====================================================================
\chapter{La interfaz}

La pantalla se divide en cuatro zonas:

\begin{description}
  \item[Barra superior.] Logo, las cuatro pestañas
    (\pestana{Explorar}, \pestana{Comparar}, \pestana{Umbrales} e
    \pestana{Instrucciones}) y, a la derecha, el selector
    \campo{Umbrales}, que indica qué perfil se aplica a gráficos y alertas.
  \item[Panel lateral «Archivos y series».] Zona para cargar archivos, lista
    de archivos cargados, buscador y lista de series para marcar.
  \item[Área principal.] Muestra el contenido de la pestaña activa.
  \item[Pie de página.] Autor y versión.
\end{description}

\section{Tema claro y oscuro}

La aplicación sigue automáticamente la preferencia de colores del sistema
operativo. Los gráficos exportados usan los colores del tema activo: si se
exporta en modo oscuro, la imagen tendrá fondo oscuro.

\section{Pantallas pequeñas}

En pantallas de hasta unos 820\,px de ancho (teléfonos y tabletas):

\begin{itemize}
  \item El panel lateral se oculta y se abre con el botón
        \boton{Archivos} de la barra superior.
  \item Las tablas se transforman en \emph{tarjetas}: cada fila se muestra como
        un bloque con el nombre de cada columna sobre su valor.
\end{itemize}

% =====================================================================
\chapter{Cargar archivos}

\section{Cómo cargar}

\begin{itemize}
  \item \textbf{Arrastrar y soltar} uno o varios archivos en cualquier parte de
        la ventana, o
  \item hacer clic en el recuadro \emph{«Arrastra tus CSV aquí»} y elegirlos
        (se aceptan \texttt{.csv}, \texttt{.tsv} y \texttt{.txt}).
\end{itemize}

Se pueden cargar varios archivos a la vez y seguir agregando más. Cada
archivo aparece en el panel lateral con:

\begin{itemize}
  \item su nombre;
  \item el \textbf{origen}: el servidor (y la instancia SQL con nombre, si la
        hay, por ejemplo \texttt{SRV01\bs VENTAS});
  \item el número de filas y el rango de fechas.
\end{itemize}

El botón \boton{\texttimes} a la derecha quita el archivo.

\section{Qué detecta la aplicación automáticamente}

\subsection{Formato del archivo}
\begin{itemize}
  \item \textbf{Separador:} coma (\texttt{,}), punto y coma (\texttt{;}),
        tabulación o barra vertical (\texttt{|}). Se elige el que más aparece
        en la primera línea, sin contar los que están entre comillas.
  \item Codificación con o sin BOM, saltos de línea Windows o Unix, campos
        entre comillas y comillas escapadas (\texttt{""}).
  \item \textbf{Coma decimal:} si el separador es punto y coma, se asume que los
        números usan coma decimal (\texttt{1,5} = 1.5), como en las
        configuraciones regionales en español.
\end{itemize}

\subsection{Columna de fecha y hora}
Se usa la \textbf{primera columna} en la que al menos el 95\,\% de los primeros
200 valores no vacíos parezcan una fecha. Formatos admitidos:

\begin{center}
\begin{tabular}{ll}
  \toprule
  Formato & Ejemplo \\
  \midrule
  Año-mes-día & \texttt{2026-07-25 14:30:00} \\
  Día/mes/año & \texttt{25/07/2026 14:30} \\
  Mes/día/año (Perfmon en inglés) & \texttt{07/25/2026 14:30:00.500} \\
  \bottomrule
\end{tabular}
\end{center}

\begin{itemize}
  \item Separadores de fecha: \texttt{/}, \texttt{-} o \texttt{.}; la hora
        puede ir tras un espacio o una \texttt{T}; los segundos y las
        fracciones de segundo son opcionales.
  \item Años de dos dígitos se interpretan como 20\emph{xx}.
  \item \textbf{Día/mes o mes/día:} si algún valor del primer campo es mayor que
        12, el formato es día/mes/año. Si ninguno de los dos campos supera 12,
        no se puede saber, y se asume \textbf{mes/día/año} (lo que usa
        Perfmon), con un aviso.
  \item Las horas se muestran \textbf{tal como vienen en el archivo}, sin
        conversión de zona horaria.
\end{itemize}

\subsection{Columnas de datos}
\begin{itemize}
  \item Si al menos el 90\,\% de los valores no vacíos de una columna son
        números, es una \textbf{serie} (se puede graficar).
  \item Si no, es una columna de texto (\emph{etiqueta}); no se grafica, pero se
        usa para identificar el origen cuando no hay cabecera de Perfmon.
  \item Las celdas vacías o con solo espacios se tratan como «sin dato» y
        aparecen como huecos en el gráfico.
\end{itemize}

\subsection{Cabeceras de Perfmon}
Las columnas con formato
\texttt{\bs\bs SERVIDOR\bs Objeto\bs Contador} se reconocen y se
muestran como \emph{Objeto · Contador}. Además:
\begin{itemize}
  \item Se quita el prefijo de instancia de SQL Server
        (\texttt{SQLServer:} para la instancia por defecto o
        \texttt{MSSQL\$NOMBRE:} para una instancia con nombre). Así, el mismo
        contador de dos servidores o instancias distintas se considera
        \textbf{la misma métrica} y se puede superponer y comparar.
  \item Se deduce la unidad: \% (contadores que empiezan con \% o contienen
        \emph{ratio}), /s (terminan en \texttt{/sec}), KB, MB o ms.
  \item El servidor de origen del archivo es el que más se repite en las
        cabeceras.
\end{itemize}

\subsection{Calidad del dato}
Al cargar, la aplicación también:
\begin{itemize}
  \item ordena por fecha las filas que vengan desordenadas;
  \item descarta las filas con fecha repetida (se conserva la primera);
  \item descarta las filas con fecha inválida.
\end{itemize}
Cada una de estas correcciones genera un aviso en color ámbar en el panel
lateral, bajo el área de carga.

\section{Mensajes al cargar}

\begin{center}
\begin{tabularx}{\textwidth}{L{6.2cm}X}
  \toprule
  Mensaje & Significado y solución \\
  \midrule
  \emph{No se encontró una columna de fecha/hora reconocible} &
    Ninguna columna tiene fechas en un formato admitido. Revisa que el archivo
    tenga una columna de fecha y que el separador sea el correcto. \\
  \emph{El archivo no tiene filas de datos} &
    El archivo está vacío o solo tiene la cabecera. \\
  \emph{sin filas con fecha válida} &
    Se encontró la columna de fecha, pero ninguna fila se pudo interpretar. \\
  \emph{Formato de fecha ambiguo: se asumió mes/día/año.} &
    Aviso: comprueba en el gráfico que las fechas sean correctas. \\
  \emph{N filas venían fuera de orden y se ordenaron por fecha.} &
    Aviso informativo; los datos ya están corregidos. \\
  \emph{N filas con fecha repetida se descartaron.} &
    Aviso informativo. \\
  \emph{N filas con fecha inválida se descartaron.} &
    Aviso informativo. \\
  \bottomrule
\end{tabularx}
\end{center}

\section{Seleccionar series}

Debajo de los archivos aparece la lista de series agrupada por origen.

\begin{itemize}
  \item Usa \campo{Buscar serie\ldots} para filtrar por nombre.
  \item Marca la casilla de cada serie que quieras ver. Cada serie marcada
        recibe un color, que se muestra con un punto junto a su nombre.
  \item Se pueden marcar hasta \textbf{8 series} a la vez.
\end{itemize}

% =====================================================================
\chapter{Pestaña Explorar}

Es la vista principal. Muestra un gráfico por cada \textbf{contador}
seleccionado y, debajo de cada gráfico, su registro de alertas.

\begin{consejo}
Si marcas el mismo contador (por ejemplo \emph{\% Processor Time}) en
archivos de varios servidores, todos se \textbf{superponen en un solo
gráfico}, con un color por servidor. Así se comparan servidores de un
vistazo.
\end{consejo}

\section{Partes del gráfico}

\begin{description}
  \item[Título y subtítulo.] Nombre del contador y, si el perfil activo tiene
    umbral para él, su \emph{valor ideal} (por ejemplo «Valor ideal: $<$ 80\,\%»).
  \item[Leyenda.] Un elemento por servidor. Si un mismo servidor aparece en
    dos archivos, se añade el nombre del archivo entre paréntesis.
  \item[Eje vertical.] Escala con valores redondos. Si todas las series son
    porcentajes, el eje no pasa de 100\,\%.
  \item[Eje horizontal.] Se adapta a la duración visible: horas si es menos de
    dos días, mes-día y hora si es menos de veinte días, y fechas completas
    para periodos más largos.
  \item[Franjas de umbral.] Zona \warnz{} en ámbar y zona \critz{} en rojo, con
    líneas discontinuas en cada límite y una etiqueta como
    «Warning Zone $>$ 80».
  \item[Huecos sin datos.] Cuando falta un tramo largo de datos (más de 12
    veces el intervalo normal de muestreo), los extremos se unen con una línea
    punteada tenue para que el hueco sea visible.
  \item[Tabla «Estadísticas generales».] Al pie del gráfico, una fila por
    serie con:
    \begin{itemize}
      \item \campo{Condición}: \okz{} (nunca salió del rango ideal), \warnz{} o
            \critz{} (la peor alerta detectada, con el límite y el número de
            eventos) o \textsf{SIN UMBRAL}.
      \item \campo{Mínimo}, \campo{Promedio} y \campo{Máximo} de toda la serie
            cargada (no solo del tramo visible con zoom).
    \end{itemize}
    Esta tabla forma parte de la imagen, así que se incluye al exportar.
\end{description}

\section{Interacción con el gráfico}

\begin{itemize}
  \item \textbf{Leer valores:} al pasar el cursor aparece una línea vertical y
        un recuadro con la fecha y el valor de cada serie en ese instante.
  \item \textbf{Acercar (zoom):} haz clic y arrastra horizontalmente sobre el
        tramo de interés.
  \item \textbf{Volver a la vista completa:} doble clic sobre el gráfico o
        botón \boton{Restablecer zoom}.
\end{itemize}

El zoom se aplica a \textbf{todos} los gráficos de la pestaña a la vez, para
poder comparar el mismo intervalo en distintos contadores.

\section{Alertas detectadas}

Debajo de cada gráfico con umbral aparece la sección
\textbf{Alertas detectadas}. Una \emph{alerta} (o \emph{evento}) es un tramo
de muestras \textbf{consecutivas} fuera del rango ideal. El tramo se corta
cuando el valor vuelve al rango ideal o cuando falta un tramo largo de datos.

\begin{center}
\begin{tabularx}{\textwidth}{lX}
  \toprule
  Columna & Contenido \\
  \midrule
  \campo{Fecha / hora} & Inicio del evento. \\
  \campo{Severidad} & \critz{} si en algún momento entró en la zona crítica;
                      si no, \warnz. \\
  \campo{Servidor} & Solo si hay varios servidores en el gráfico. \\
  \campo{Problema} & Contador y límite superado, por ejemplo
                     «\% Processor Time $>$ 90». \\
  \campo{Valor pico} & Valor más extremo alcanzado durante el evento. \\
  \campo{Duración} & Desde la primera hasta la última muestra del evento
                     («1 muestra» si fue un solo punto). \\
  \campo{Acción} & Botón \boton{Ver en gráfico}. \\
  \bottomrule
\end{tabularx}
\end{center}

Controles de la sección:
\begin{itemize}
  \item Contadores de \critz{} y \warnz{} totales.
  \item Filtro: \emph{Todas}, \emph{Critical} o \emph{Warning}.
  \item Orden: \emph{Por fecha} o \emph{Críticas primero}.
  \item Se muestran 10 eventos; \boton{Mostrar más} añade 20 cada vez.
\end{itemize}

\boton{Ver en gráfico} acerca todos los gráficos al periodo del evento (con un
margen alrededor), lo marca con una franja gris y señala el valor pico con un
punto. Si la serie no estaba marcada, se marca sola. Para quitar la marca usa
\boton{Quitar resaltado}.

Si el contador no tiene umbral en el perfil activo (o es «Solo informativo»),
se indica «Sin umbral para este contador: no genera alertas». Si tiene umbral
y nunca salió del rango, se indica «Sin alertas».

\section{Exportar gráficos}

En la barra superior de la pestaña:

\begin{itemize}
  \item Selector de formato: \textbf{PNG} (recomendado para informes, nítido) o
        \textbf{JPEG} (archivo más liviano).
  \item \boton{Descargar gráfico}: guarda la imagen. Si hay varios gráficos, el
        botón se llama \boton{Descargar todos} y los apila en
        \textbf{una sola imagen}.
  \item \boton{Copiar al portapapeles} / \boton{Copiar todos}: copia la imagen
        (siempre en PNG) para pegarla directamente en Word, PowerPoint, correo,
        etc.
\end{itemize}

Cuando hay más de un gráfico, al pasar el cursor sobre cada uno aparecen
\boton{Copiar} y \boton{Descargar} para exportarlo individualmente. El archivo
se llama \texttt{grafico-AAAA-MM-DD.png} (o \texttt{.jpg}) con la fecha del
día.

La imagen exportada incluye título, leyenda, franjas, gráfico y tabla de
estadísticas, en alta resolución.

% =====================================================================
\chapter{Pestaña Comparar}

Permite superponer \textbf{dos periodos} de una misma métrica para ver cómo
cambió: antes y después de un cambio, mes contra mes, servidor contra
servidor, etc.

\section{Configurar la comparación}

\begin{enumerate}
  \item En \textbf{Periodo A} y \textbf{Periodo B} elige el archivo y las
        fechas \emph{desde} y \emph{hasta} (ambos días se incluyen completos).
  \item En \campo{Métrica} elige el contador. Solo aparecen los que existen en
        ambos periodos.
  \item En \campo{Eje X} elige el modo:
    \begin{description}
      \item[Superpuesto.] Ambos periodos empiezan en el mismo punto (día 0)
        y se dibujan uno encima del otro. Ideal para comparar periodos de
        fechas diferentes. Al pasar el cursor se muestra la fecha real de cada
        periodo.
      \item[Calendario real.] Cada periodo se dibuja en sus fechas reales.
    \end{description}
\end{enumerate}

\begin{nota}[Valores iniciales]
Con un solo archivo cargado, la aplicación propone dividir su rango en dos
mitades (A = primera mitad, B = segunda). Con dos o más archivos, propone el
primero como A y el segundo como B, cada uno con su rango completo. Al
cambiar el archivo de un periodo, las fechas se ajustan al rango completo de
ese archivo.
\end{nota}

\section{Resultado}

Se muestra un gráfico con ambos periodos (A en azul, B en rojo), las franjas
de umbral si corresponde, y la tabla de estadísticas con tres filas:

\begin{itemize}
  \item \textbf{Periodo A} y \textbf{Periodo B}: condición, mínimo, promedio y
        máximo \emph{del periodo elegido}.
  \item \textbf{Cambio B vs A}: variación porcentual de mínimo, promedio y
        máximo:
        \[
          \Delta\% = \frac{B - A}{|A|} \times 100
        \]
        En azul los aumentos y en ámbar las disminuciones. Si el valor de A es
        0, se muestra «–».
\end{itemize}

El gráfico se exporta igual que en \pestana{Explorar}. En esta pestaña no hay
zoom con el ratón; para acotar, ajusta las fechas de los periodos.

Si aparece «No hay datos comparables», revisa que las fechas tengan datos en
ambos periodos y que compartan al menos una métrica.

% =====================================================================
\chapter{Umbrales}

\section{Conceptos}

\begin{description}
  \item[Perfil de umbrales.] Conjunto de reglas de una empresa o cliente
    (por ejemplo \emph{SII} o \emph{TGR}). Solo hay un perfil activo a la vez;
    se elige en el selector \campo{Umbrales} de la barra superior o en la
    pestaña \pestana{Umbrales}. Al cambiarlo, los gráficos y alertas se
    recalculan al instante.
  \item[Regla.] Fila del perfil que dice qué contador vigilar y con qué
    límites. Campos:
    \begin{itemize}
      \item \campo{Contador}: nombre descriptivo.
      \item \campo{Patrón}: texto que debe aparecer en el nombre del contador
            del CSV (ver más abajo).
      \item \campo{Valor ideal}: texto libre que se muestra en el gráfico
            (por ejemplo «$<$ 80\,\%»). Si se deja vacío, se genera solo.
      \item \campo{Criterio}: cuándo se activa la alerta.
      \item \campo{Warning} y \campo{Critical}: valores donde empiezan la zona
            ámbar y la roja.
    \end{itemize}
\end{description}

\subsection{Criterio y zonas}

\begin{center}
\begin{tabularx}{\textwidth}{L{3.6cm}L{2.6cm}X}
  \toprule
  Criterio & Ejemplo & Zonas (\emph{v} = valor medido) \\
  \midrule
  \emph{Alerta si sube de ($>$)} \newline (menor es mejor) &
    CPU $<$ 80\,\% &
    \okz: $v \le$ Warning \newline
    \warnz: Warning $< v \le$ Critical \newline
    \critz: $v >$ Critical \\
  \addlinespace
  \emph{Alerta si baja de ($<$)} \newline (mayor es mejor) &
    Cache hit $>$ 95\,\% &
    \okz: $v \ge$ Warning \newline
    \warnz: Critical $\le v <$ Warning \newline
    \critz: $v <$ Critical \\
  \addlinespace
  \emph{Solo informativo} &
    Transactions/sec &
    No genera alertas ni franjas (valores «Variable»). \\
  \bottomrule
\end{tabularx}
\end{center}

Por lo tanto, con \emph{sube de}, Critical debe ser \textbf{mayor} que
Warning; con \emph{baja de}, \textbf{menor}. Las tablas de referencia de los
clientes suelen indicar solo el valor ideal: úsalo como Warning y define tú
un margen razonable para Critical.

\subsection{Patrones}

\begin{itemize}
  \item \textbf{Texto simple:} la regla se aplica a toda serie cuyo nombre
        \emph{contenga} ese texto, sin distinguir mayúsculas. Ejemplo:
        \texttt{Lock Requests/sec}.
  \item \textbf{Expresión regular:} si el patrón empieza y termina con
        \texttt{/}, se interpreta como expresión regular (sin distinguir
        mayúsculas). Ejemplo: \texttt{/Processor.*Time\$/}.
  \item El patrón se compara con la clave \emph{Objeto\bs Contador} (sin
        servidor ni prefijo de instancia SQL) o, en CSV que no son de Perfmon,
        con el nombre de la columna.
  \item \textbf{Prioridad:} si varias reglas coinciden con una serie, se aplica
        \textbf{la primera de la lista}. Usa las flechas para reordenar.
\end{itemize}

\subsection{Variantes (perfil de uso)}
Algunas reglas tienen valores distintos según el tipo de servidor. Por
ejemplo, en el perfil SII, \emph{Buffer cache hit ratio} tiene las variantes
\textbf{OLTP} ($>$ 95\,\%) y \textbf{OLAP} ($>$ 80\,\%). La variante se elige
en la propia fila; al cambiarla, Warning y Critical toman los valores de esa
variante, y si luego los editas, se guardan en la variante activa.

\subsection{Reglas de relación}
La regla \emph{Target vs Total Memory} del perfil SII no compara contra un
número fijo: mide el \textbf{porcentaje del tiempo en que Total Server Memory
supera a Target Server Memory}. No dibuja franjas ni genera eventos; su
resultado aparece en la tabla \emph{Cumplimiento por archivo}.

\section{Perfiles de fábrica}

\subsection*{SII}
\begin{center}
\small
\begin{tabularx}{\textwidth}{XL{3.4cm}lrr}
  \toprule
  Contador & Valor ideal & Criterio & Warning & Critical \\
  \midrule
  Memory\bs Pages/sec & $<$ 20 & sube de & 20 & 40 \\
  \% Processor Time & $<$ 80\,\% & sube de & 80 & 90 \\
  Buffer cache hit ratio & $>$ 95\,\% (OLTP) \newline $>$ 80\,\% (OLAP) & baja de &
    95 / 80 & 90 / 70 \\
  Transactions/sec & Variable & informativo & -- & -- \\
  Sesiones por día (User Connections) & Variable & informativo & -- & -- \\
  SQL Compilations/sec & $\le$ 100 & sube de & 100 & 200 \\
  Target vs Total Memory & Total $\le$ Target & relación & -- & -- \\
  Locks: Lock Requests/sec & $<$ 1000 & sube de & 1000 & 2000 \\
  \bottomrule
\end{tabularx}
\end{center}

\subsection*{TGR}
\begin{center}
\small
\begin{tabularx}{\textwidth}{XL{3.4cm}lrr}
  \toprule
  Contador & Valor ideal & Criterio & Warning & Critical \\
  \midrule
  \% Processor Time & $<$ 80\,\% & sube de & 80 & 90 \\
  Cache Hit Ratio & $>$ 80\,\% & baja de & 80 & 70 \\
  User Connections & Variable & informativo & -- & -- \\
  Costo de CPU por BDD & Variable & informativo & -- & -- \\
  Gráfica Crecimiento BDD & Variable & informativo & -- & -- \\
  \bottomrule
\end{tabularx}
\end{center}

En ambos perfiles, el valor \emph{Warning} es el valor ideal de la tabla de
referencia del cliente y el valor \emph{Critical} es un margen propuesto, que
se puede editar.

\section{Pestaña Umbrales: editar perfiles}

\subsection{Barra de perfiles}
\begin{itemize}
  \item Selector del perfil activo.
  \item \boton{+ Nuevo}: crea un perfil vacío y abre la edición de sus
        detalles.
  \item \boton{Exportar}: descarga el perfil activo como
        \texttt{umbrales-<nombre>.json}.
  \item \boton{Importar}: carga un archivo \texttt{.json} de perfiles (ver
        capítulo~\ref{cap:json}). Los perfiles importados se \textbf{agregan}
        (no reemplazan a los existentes); si el nombre ya existe se le añade
        un número, como «Cliente (2)».
  \item Menú \boton{$\cdots$}:
    \begin{itemize}
      \item \emph{Duplicar perfil}: crea una copia editable.
      \item \emph{Exportar todos los perfiles}: un solo archivo
            \texttt{umbrales-todos.json}.
      \item \emph{Restablecer valores de origen}: solo en SII y TGR; vuelve a
            los valores de fábrica de ese perfil.
      \item \emph{Eliminar perfil\ldots}: solo en perfiles creados por el
            usuario; pide confirmación.
      \item \emph{Restablecer perfiles de fábrica}: deja únicamente SII y TGR
            con sus valores originales; \textbf{elimina todos los demás
            perfiles y cambios}. Pide confirmación.
    \end{itemize}
  \item Indicador \textcolor{warnc}{\textbf{$\bullet$ Cambios sin exportar}}:
        el perfil tiene cambios que aún no se guardaron en un archivo.
\end{itemize}

\subsection{Detalles del perfil}
Con \boton{Editar detalles} se cambian el \textbf{nombre}, la
\textbf{descripción} (de dónde salen los valores) y la marca
\textbf{Borrador}. Un perfil en borrador muestra un aviso: sus alertas usan
valores aún no confirmados.

\subsection{Tabla de reglas}
Cada fila es una regla y se edita directamente en la tabla; los cambios se
guardan y se aplican al instante.

\begin{itemize}
  \item Flechas \boton{$\blacktriangle$} \boton{$\blacktriangledown$}:
        cambian la prioridad.
  \item Bajo el nombre se edita el patrón (en letra monoespaciada).
  \item Con archivos cargados, una etiqueta indica cuántas series usan la
        regla, o \emph{«sin coincidencias»} en ámbar si el patrón no coincide
        con ninguna. Es la forma más rápida de comprobar un patrón.
  \item \boton{\texttimes}: elimina la regla.
  \item \boton{+ Agregar contador}: añade una regla nueva al final.
\end{itemize}

Los campos con problemas se marcan en rojo, con la explicación al pasar el
cursor:
\begin{itemize}
  \item \emph{Falta el patrón: sin él la regla no coincide con ninguna serie.}
  \item \emph{La expresión regular del patrón no es válida.}
  \item \emph{Con «menor es mejor», Critical debe ser mayor que Warning.}
  \item \emph{Con «mayor es mejor», Critical debe ser menor que Warning.}
\end{itemize}

\subsection{Cumplimiento por archivo}
Al final de la pestaña, para cada archivo cargado y cada serie con regla, se
muestra el \textbf{porcentaje de muestras} que estuvo en zona Warning y en
zona Critical. Para la regla de relación de memoria se muestra el porcentaje
del tiempo con Total $>$ Target.

\section{Dónde se guardan los perfiles}

Los perfiles se guardan automáticamente en el navegador. Esto significa que:

\begin{aviso}
\begin{itemize}[leftmargin=*]
  \item Son propios de \textbf{ese navegador en ese equipo}: no se comparten
        con otros equipos ni otros navegadores.
  \item Se pierden si se borran los datos de navegación o se usa una ventana
        privada o de incógnito.
\end{itemize}
Por eso conviene \textbf{exportar} los perfiles y guardar los \texttt{.json}
en una carpeta compartida o en el repositorio. Para pasarlos a otro equipo,
basta con importarlos allí.
\end{aviso}

Cuando una nueva versión de la aplicación cambia los valores de fábrica de SII
o TGR, los perfiles que no editaste se actualizan solos. Si los editaste, se
conservan tus cambios y aparece un aviso para que, si quieres, uses
\emph{Restablecer valores de origen}.

% =====================================================================
\chapter{Formato de archivo de perfiles (JSON)}
\label{cap:json}

\section{Estructura}

Un perfil es un archivo JSON con esta forma (es la plantilla estándar, que
también está en \texttt{perfiles/plantilla.json} y en la pestaña
\pestana{Instrucciones}):

\begin{lstlisting}
{
  "format": "analizador-csv.umbrales",
  "version": 1,
  "name": "Nombre del perfil (empresa o cliente)",
  "description": "De dónde salen estos valores (documento, fecha, responsable…)",
  "draft": false,
  "rules": [
    {
      "id": "cpu",
      "name": "% Processor Time",
      "ideal": "< 80 %",
      "pat": "Processor(_Total)\\% Processor Time",
      "dir": "lt",
      "warn": 80,
      "crit": 90
    },
    {
      "id": "cache",
      "name": "Buffer cache hit ratio",
      "ideal": "> 95 %",
      "pat": "Buffer cache hit ratio",
      "dir": "gt",
      "warn": 95,
      "crit": 90
    }
  ]
}
\end{lstlisting}

\section{Campos}

\begin{center}
\small
\begin{tabularx}{\textwidth}{lX}
  \toprule
  Campo & Descripción \\
  \midrule
  \texttt{format} & Siempre \texttt{"analizador-csv.umbrales"}. \\
  \texttt{version} & Siempre \texttt{1}. \\
  \texttt{name} & Nombre del perfil (empresa o cliente). \\
  \texttt{description} & Opcional. Origen de los valores. \\
  \texttt{draft} & Opcional. \texttt{true} si es borrador. \\
  \texttt{updated} & Opcional. Fecha de exportación (la añade la aplicación). \\
  \texttt{rules} & Lista de reglas. \\
  \midrule
  \texttt{rules[].id} & Identificador corto, único dentro del perfil. \\
  \texttt{rules[].name} & Nombre del contador que se muestra. \\
  \texttt{rules[].ideal} & Texto del valor ideal. \\
  \texttt{rules[].pat} & Patrón (texto o \texttt{/regex/}). En JSON la barra
                         invertida se escribe doble: \texttt{\bs\bs}. \\
  \texttt{rules[].dir} & \texttt{"lt"} (alerta si sube), \texttt{"gt"}
                         (alerta si baja) o \texttt{"none"} (informativo). \\
  \texttt{rules[].warn}, \texttt{crit} & Límites numéricos. Con
                         \texttt{"none"} usar 0. \\
  \texttt{rules[].variants} & Opcional. Valores alternativos, p. ej.
     \texttt{\{"OLTP": \{"warn": 95, "crit": 90\}, "OLAP": \{\ldots\}\}}. \\
  \texttt{rules[].variant} & Opcional. Variante activa. \\
  \texttt{rules[].relation} & Opcional. \texttt{true} para la regla de
                              relación de memoria. \\
  \bottomrule
\end{tabularx}
\end{center}

Un archivo también puede contener \textbf{varios perfiles}: es lo que genera
\emph{Exportar todos los perfiles}
(\texttt{\{"format": \ldots, "version": 1, "profiles": [ \ldots ]\}}). Se
importa del mismo modo.

\section{Validación al importar}

La aplicación corrige automáticamente los detalles menores: un criterio
desconocido pasa a \texttt{"none"}, un número inválido pasa a 0, un
\texttt{id} vacío o repetido se reemplaza por uno único y una regla sin nombre
recibe «Contador \emph{N}». Los errores graves se informan con un mensaje:

\begin{itemize}
  \item \emph{El archivo no es un JSON válido.}
  \item \emph{El archivo no es un perfil de umbrales de esta herramienta.}
  \item \emph{El archivo no tiene la lista de reglas («rules»).}
  \item \emph{El archivo no contiene perfiles.}
\end{itemize}

% =====================================================================
\chapter{Pestaña Instrucciones y perfiles asistidos por IA}

La pestaña \pestana{Instrucciones} reúne una guía rápida, la explicación de
los umbrales, la plantilla JSON (con los botones \boton{Copiar plantilla} y
\boton{Descargar plantilla (.json)}) y un \textbf{prompt} listo para usar con
un asistente de IA.

\section{Crear el perfil de un cliente nuevo con IA}

Cuando un cliente entrega su tabla de valores ideales (en texto, PDF o
imagen):

\begin{enumerate}
  \item En \pestana{Instrucciones}, pulsa \boton{Copiar prompt}.
  \item Pégalo en el asistente de IA de tu preferencia y adjunta o pega la
        tabla del cliente donde indica
        \emph{[PEGA AQUÍ LA TABLA O ADJUNTA LA IMAGEN]}.
  \item Guarda la respuesta (solo el bloque JSON) en un archivo de texto con
        extensión \texttt{.json}.
  \item En \pestana{Umbrales}, pulsa \boton{Importar} y elige el archivo.
  \item Revisa la tabla de reglas: corrige lo marcado en rojo, verifica que
        cada regla tenga series asociadas y ajusta los valores
        \emph{Critical} que la IA haya propuesto por su cuenta (el prompt le
        pide que los indique al final de su respuesta).
  \item Si los valores aún no están confirmados por el cliente, marca el
        perfil como \textbf{Borrador}.
  \item Exporta el perfil para guardar una copia.
\end{enumerate}

\begin{aviso}
La IA puede equivocarse al transcribir valores o elegir el criterio. Revisa
siempre el perfil generado contra la tabla original antes de usarlo en un
informe.
\end{aviso}

% =====================================================================
\chapter{Flujos de trabajo habituales}

\section{Informe de rendimiento de un servidor}
\begin{enumerate}
  \item Carga el CSV de Perfmon del servidor.
  \item Elige el perfil del cliente en el selector \campo{Umbrales}.
  \item En el buscador, filtra y marca los contadores del informe.
  \item Revisa la condición de cada contador en su tabla de estadísticas y
        las alertas debajo de cada gráfico; usa \boton{Ver en gráfico} para
        analizar los eventos críticos.
  \item Con \boton{Restablecer zoom} vuelve a la vista completa y exporta con
        \boton{Descargar todos} o copia cada gráfico al informe.
  \item En \pestana{Umbrales}, consulta la tabla \emph{Cumplimiento por
        archivo} para obtener los porcentajes de tiempo en Warning y Critical.
\end{enumerate}

\section{Comparar varios servidores}
Carga un CSV por servidor y marca el mismo contador en cada uno: en
\pestana{Explorar} se superponen en un único gráfico, con una fila por
servidor en la tabla de estadísticas.

\section{Antes y después de un cambio}
Carga el archivo (o los dos archivos) y en \pestana{Comparar} define el
periodo A (antes) y el B (después), en modo \emph{Superpuesto}. La fila
\emph{Cambio B vs A} da la variación porcentual.

% =====================================================================
\chapter{Solución de problemas}

\begin{description}
  \item[No aparece mi archivo / «no se pudo leer».] Revisa el mensaje en el
    panel lateral (ver la tabla de mensajes del capítulo de carga). Lo más
    común es que la columna de fecha tenga un formato no admitido.
  \item[Las fechas aparecen con día y mes invertidos.] Ocurre si todas las
    fechas tienen día $\le$ 12 y no son del tipo año-mes-día; la aplicación
    asume mes/día/año y lo avisa. Si es posible, exporta el CSV con fechas
    \texttt{AAAA-MM-DD}.
  \item[Los números salen mal o la columna no aparece como serie.] Si el CSV
    usa coma decimal, el separador de columnas debe ser punto y coma. Una
    columna con más de un 10\,\% de valores no numéricos se trata como texto.
  \item[No puedo marcar más series.] El máximo es 8 series simultáneas.
    Desmarca alguna.
  \item[Un contador no muestra franjas ni alertas.] El perfil activo no tiene
    regla para él, la regla es «Solo informativo», o el patrón no coincide.
    En \pestana{Umbrales}, mira si la regla dice \emph{«sin coincidencias»}.
  \item[Se aplica una regla distinta a la esperada.] Hay otra regla más arriba
    cuyo patrón también coincide. Haz el patrón más específico o sube la
    regla correcta.
  \item[«El navegador no permitió copiar».] Usa \boton{Descargar gráfico}, o
    usa la versión en línea (\texttt{https}).
  \item[Perdí mis perfiles.] Se borraron los datos del navegador o se usó otro
    navegador o equipo. Importa la última copia exportada.
  \item[La imagen exportada tiene fondo oscuro.] El sistema está en modo
    oscuro; cámbialo a modo claro antes de exportar si el informe lo
    requiere.
\end{description}

% =====================================================================
\appendix

\chapter{Información para desarrolladores}

\section{Estructura del proyecto}

\begin{center}
\small
\begin{tabularx}{\textwidth}{lX}
  \toprule
  Ruta & Contenido \\
  \midrule
  \texttt{src/core/parser/} & Lectura de CSV (\texttt{csv.ts}), fechas
     (\texttt{time.ts}), cabeceras Perfmon (\texttt{perfmon.ts}) y armado del
     conjunto de datos y estadísticas (\texttt{dataset.ts}). TypeScript. \\
  \texttt{src/core/thresholds.ts} & Reglas, zonas, alertas, perfiles de
     fábrica, importación y exportación de perfiles, plantilla y prompt. \\
  \texttt{src/core/compare.ts} & Recorte por fechas y comparación de periodos. \\
  \texttt{src/ui/app.js} & Interfaz: JavaScript sin frameworks y gráficos en
     \emph{canvas}. \\
  \texttt{src/ui/index.template.html} & Plantilla HTML y estilos. \\
  \texttt{src/assets/} & Logos y favicon. \\
  \texttt{perfiles/} & Perfiles SII y TGR y plantilla, como \texttt{.json}
     descargables. \\
  \texttt{scripts/build.mjs} & Une todo en \texttt{dist/index.html}. \\
  \texttt{tests/} & Pruebas del núcleo. \\
  \texttt{.github/workflows/deploy.yml} & Publicación en GitHub Pages. \\
  \bottomrule
\end{tabularx}
\end{center}

\section{Comandos}

Requiere Node.js 22 o superior (el build usa
\texttt{stripTypeScriptTypes} de \texttt{node:module}).

\begin{lstlisting}
npm install         # dependencias de desarrollo (TypeScript)
npm run build       # genera dist/index.html (un solo archivo)
npm test            # pruebas del núcleo
npm run typecheck   # verificación de tipos
\end{lstlisting}

\section{Publicación}

Cada \emph{push} a la rama \texttt{main} ejecuta el flujo de GitHub Actions,
que compila el proyecto y publica \texttt{dist/} en GitHub Pages. La versión
del pie de página se toma de \texttt{package.json} (\texttt{1.0.0} se muestra
como «Versión 1.0»).

\section{Mantener los perfiles de fábrica}

Los perfiles SII y TGR están definidos en \texttt{src/core/thresholds.ts}
(\texttt{SII\_RULES}, \texttt{TGR\_RULES} y \texttt{BUILTIN\_PROFILES}) y
copiados en \texttt{perfiles/*.json}. Al cambiar sus valores:
\begin{enumerate}
  \item edita \texttt{thresholds.ts};
  \item sube el número \texttt{rev} del perfil en \texttt{BUILTIN\_PROFILES}
        (así los navegadores de los usuarios que no lo editaron se actualizan
        solos);
  \item actualiza el \texttt{.json} correspondiente en \texttt{perfiles/}. La
        prueba \texttt{tests/perfiles.test.ts} falla si quedan distintos.
\end{enumerate}

% =====================================================================
\chapter{Glosario}

\begin{description}
  \item[Contador.] Métrica medida por Perfmon, por ejemplo
    \emph{\% Processor Time}.
  \item[Serie.] Columna numérica de un CSV: los valores de un contador a lo
    largo del tiempo.
  \item[Perfmon.] \emph{Performance Monitor}, la herramienta de Windows que
    registra contadores de rendimiento y puede exportarlos a CSV.
  \item[Perfil de umbrales.] Conjunto de reglas de alerta de una empresa o
    cliente.
  \item[Warning / Critical.] Zonas de advertencia (ámbar) y crítica (rojo).
  \item[Evento o alerta.] Tramo continuo de muestras fuera del rango ideal.
  \item[OLTP / OLAP.] Tipos de carga de una base de datos: transaccional
    (muchas operaciones cortas) o analítica (consultas grandes).
  \item[Patrón.] Texto o expresión regular que asocia una regla con las
    series.
\end{description}

\end{document}
